\section{CONCLUSION}
This paper examined ability-diverse collaboration in video creation workshops for blind and visually impaired communities, revealing fundamental challenges to prevailing models of inclusive education and accessible technology design. Through detailed analysis of seven workshops combining in-person and remote delivery, we identified critical barriers and collaborative patterns that emerge when participants with different visual abilities learn video production together.

Our investigation addressed two primary research questions. First, we documented the challenges and structural barriers affecting video creation education in ability-diverse teams, identifying three interlocking asymmetries: expertise gaps between video and accessibility pedagogies, tool ecosystem fragmentation, and fundamental incompatibilities between visual and non-visual interaction paradigms. These asymmetries manifested not as simple accessibility failures but as systemic misalignments between different epistemic frameworks for understanding and interacting with video creation tools. Second, we identified four distinct collaborative patterns: haptic mediation, cross-modal translation, task delegation, and peer assistance, that participants developed to navigate these barriers. These patterns reveal how access emerges through dynamic, situated negotiations rather than universal design solutions.

The contributions of this work extend beyond documenting barriers to accessible video education. By analysing how participants collaboratively construct access through practice, we offer empirical evidence that challenges assumptions underlying both Universal Design for Learning and individual accommodation models. Our concept of "dynamic standard-making", the creation of temporary, localized standards between specific individuals—provides a theoretical framework for understanding how ability-diverse groups navigate fundamental epistemic differences. This framework has implications not only for video creation but for any domain where visual and non-visual ways of knowing must coexist.

Our findings suggest that future accessible media education should embrace rather than eliminate collaborative interdependence. The dyadic education model we propose, where participant-helper pairs are taught as collaborative units, acknowledges that distributed expertise represents a valid form of media literacy. Similarly, our documentation of peer teaching effectiveness, particularly in remote settings, indicates that BVI creators themselves represent an underutilized pedagogical resource. These insights point toward educational approaches that treat ability diversity as a creative resource rather than a deficit to overcome.
